% CP-VI-0108
% target_chapter: VI/49 — Basic Vanishing Results
% solution_provenance: FRESH_CANONICAL_AUTHORED_SOLUTION

\subsection*{CP-VI-0108: Vanishing on affine schemes for quasicoherent sheaves}

State the affine vanishing theorem for a quasicoherent sheaf $\mathcal F$ on $X=\operatorname{Spec}A$ and explain its significance for computations.

\medskip
\noindent\textbf{Solution.}

The theorem states
\[
H^i(X,\mathcal F)=0\qquad(i>0)
\]
for every quasicoherent $\mathcal F$ on an affine scheme $X$. Consequently, higher cohomological obstructions disappear on affine charts. This is why affine covers reduce many global calculations to Čech complexes built from intersections of affine opens, where the higher local cohomology vanishes.
